\scanpage{63}
% archive foliation: 65
a) Cat. des surfaces conformes à bord $X$,
avec comp. connexe \inserted{compacte} choisie $C$ du bord,
et une structure toroïdale dessus
rectifiable (?) pour la structure
analytique réelle, et

b) Cat. des surfaces conformes à
bord \inserted{$X'$}, avec \deleted{\unsure{le|un}} disque pointé $(D,a)$, $D$
$\subset\operatorname{Int}(X')$ marqué, tel que $\delta D$ soit rectifiable
dans $X'$ (i.e. l'application $\delta D\longrightarrow X'$
rectifiable, pour la structure analytique
naturelle de $\delta D$ comme bord de $D$).

De plus, Pour $(X',a)$ fixé, chaque
disque $D$ autour de $x$ définit un
``disque \deleted{canonique}\inserted{euclidien}'' dans $T_{X',a}$ (au sens de
la structure \inserted{vectorielle} conforme \inserted{(de $T_{X',a}$)}). Comme
\unsure{les|ces} disques forment un torseur sous
$\mathbb R^{*+}$ (par homothéties), \unsure{donc un espace}
de dim $1$, il est clair que l'application
